
\section{Method}

\subsection{Overview}

Adaptive SVD compresses each attention projection in four steps:

\begin{enumerate}
\item \textbf{Full SVD} ( 4.2): compute $U \Sigma V^\top$ in FP32.
\item \textbf{Energy-thresholded rank} ( 4.3): find smallest $r$ with $\mathcal{E}(r) \geq \tau$.
\item \textbf{Symmetric factorization} ( 4.4): set $U_r = U_{:, :r} \Sigma^{1/2}$,
  $V_r = \Sigma^{1/2} V_{:r, :}^\top$.
\item \textbf{Residual adapter} ( 4.5): attach zero-initialized $B A$ for recovery.

\end{enumerate}
\textbf{Figure 4.1} shows the pipeline. There is no calibration phase.

\begin{lstlisting}
    +------------------------------------------+
    |  Pretrained attention module A           |
    |  with projections W\_Q, W\_K, W\_V, W\_O     |
    +------------------+-----------------------+
                       |
                        
    +------------------------------------------+
    |  For each projection P:                  |
    |    Phase 1: SVD(W\_P) $\to$ U\_P  \_P V\_P^T     |
    |    Phase 2: Find r\_P via $\tau$               |
    |    Phase 3: Truncate + distribute  ^{1/2}|
    |    Phase 4: Add zero-init BA adapter     |
    +------------------+-----------------------+
                       |
                        
    +------------------------------------------+
    |  Compressed attention A' with:           |
    |    - Frozen U\_P, V\_P for each P          |
    |    - Trainable A\_P, B\_P for each P       |
    +------------------------------------------+
\end{lstlisting}

\subsection{Phase 1: Full SVD}

For each projection weight $W_P \in \mathbb{R}^{d_\text{out}^{(P)} \times d_\text{in}^{(P)}}$:

\[
W\_P = U\_P \Sigma\_P V\_P^\top
\tag{4.1}
\]

We compute this in \textbf{FP32} regardless of the model dtype. FP16 SVD can 
produce NaN for ill-conditioned matrices.

\textbf{Complexity.} $\mathcal{O}(d_\text{in} d_\text{out} \min(d_\text{in}, d_\text{out}))$. 
For $768 \times 768$, this is ~450M FLOPs, ~50 ms on CPU, ~5 ms on GPU.

\subsection{Phase 2: Energy-Thresholded Rank}

\subsubsection{Definition}

Given singular values $\sigma_1 \geq \sigma_2 \geq \dots$, define:

\[
s\_i = \sigma\_i^2, \quad \mathcal{E}(r) = \frac{\sum\_{i=1}^{r} s\_i}{\sum\_{j=1}^{\min(d\_\text{in}, d\_\text{out})} s\_j}
\tag{4.2}
\]

The rank is:

\[
r\_P = \min\{r : \mathcal{E}(r) \geq \tau\}
\tag{4.3}
\]

\subsubsection{Implementation}

\begin{lstlisting}
def \_compute\_energy\_rank(self, S, max\_rank, min\_rank=16):
    s\_squared = S ** 2
    total\_energy = s\_squared.sum()
    if total\_energy \textless{}= 1e-12:
        return max(1, min(min\_rank, max\_rank))
    
    cum\_energy = torch.cumsum(s\_squared, dim=0) / total\_energy
    idx = torch.where(cum\_energy \textgreater{}= self.energy\_threshold)[0]
    r = int(idx[0].item()) + 1 if len(idx) \textgreater{} 0 else max\_rank
    return max(1, min(max(r, min\_rank), max\_rank))
\end{lstlisting}

\textbf{Complexity.} $\mathcal{O}(\min(d_\text{in}, d_\text{out}))$ via \texttt{cumsum}.

\subsubsection{Minimum Rank Constraint}

We impose $r_P \geq r_\text{min}$ (default 16) to prevent numerical 
instability in very low-rank layers. This bounds $r_P$ from below, but for 
$\tau = 0.92$ and typical attention weights, the natural rank is far above 
$r_\text{min}$.

\subsubsection{Maximum Rank Cap}

We cap $r_P \leq \min(d_\text{in}, d_\text{out})$. In practice, $\tau < 1.0$ 
ensures $r_P < \min(d_\text{in}, d_\text{out})$, but the cap protects against 
edge cases.

\subsection{Phase 3: Symmetric Factorization}

\subsubsection{Factor Construction}

Given $r_P$, set:

\[
U\_{P, r} = U\_{P, :, :r\_P} \Sigma\_{P, :r\_P, :r\_P}^{1/2}, \quad
V\_{P, r} = \Sigma\_{P, :r\_P, :r\_P}^{1/2} V\_{P, :r\_P, :}^\top
\tag{4.4}
\]

Then $U_{P, r} V_{P, r} = W_{P, r_P}$ (the truncated SVD).

\subsubsection{Why Symmetric?}

Three reasons:

\begin{enumerate}
\item \textbf{Numerical stability.} Both factors have magnitude $\sqrt{\sigma_\text{max}}$,
  avoiding FP16 overflow.

\item \textbf{Balanced gradient flow.} During recovery, adapters interact symmetrically
  with both factors.

\item \textbf{Orthogonality.} When singular values are distinct, columns of $U_{P, r}$
  and rows of $V_{P, r}$ are orthogonal.

\end{enumerate}
\subsubsection{The Factorized Forward Pass}

The compressed projection computes:

\[
\text{Proj}\_P(x) = (x V\_{P, r}^\top) U\_{P, r}^\top + b\_P
\tag{4.5}
\]

\textbf{Parameter count:}

\begin{table}[H]
\centering
\small
\begin{tabular}{ll}
\toprule
\textbf{Configuration} & \textbf{Parameters} \\
\midrule
Original & $d_\text{in}^{(P)} d_\text{out}^{(P)}$ \\
Adaptive SVD & $r_P (d_\text{in}^{(P)} + d_\text{out}^{(P)})$ \\
\bottomrule
\end{tabular}
\end{table}

\textbf{Reduction factor:} $\frac{d_\text{in} d_\text{out}}{r_P (d_\text{in} + d_\text{out})}$.

\subsection{Phase 4: Residual Adapter}

\subsubsection{The Residual}

The truncation residual is:

\[
\Delta W\_P = W\_P - U\_{P, r} V\_{P, r}
\tag{4.6}
\]

Its Frobenius norm is $\sqrt{(1 - \mathcal{E}(r_P))} \cdot \|W_P\|_F \leq \sqrt{1 - \tau} \|W_P\|_F$. 
At $\tau = 0.92$, this is ~28\% of $\|W_P\|_F$.

\subsubsection{Adapter Parameterization}

We approximate the residual with a rank-$r_L$ factor:

\[
\Delta W\_P \approx B\_P A\_P, \quad
A\_P \in \mathbb{R}^{r\_L \times d\_\text{in}}, \quad
B\_P \in \mathbb{R}^{d\_\text{out} \times r\_L}
\tag{4.7}
\]

\subsubsection{Initialization}

\[
B\_P = 0, \quad A\_P \sim \mathcal{N}(0, 1/\sqrt{r\_L})
\tag{4.8}
\]

The zero initialization of $B_P$ ensures $\Delta W_P = B_P A_P = 0$ at step 0.

\subsubsection{Full Forward Pass}

\[
\text{Proj}\_P(x) = (x V\_{P, r}^\top) U\_{P, r}^\top + (x A\_P^\top) B\_P^\top + b\_P
\tag{4.9}
\]

\textbf{Parameter overhead:} $r_L (d_\text{in} + d_\text{out})$ per projection.

\subsection{The Full Algorithm}

\textbf{Algorithm 1:} Adaptive SVD compression for one attention module
\begin{lstlisting}
Input:  pretrained attention A with projections P $\in$ {Q, K, V, O}
        energy threshold $\tau$ $\in$ (0, 1)
        residual rank r\_L
        minimum rank r\_min
Output: compressed attention A'

1:  for P in {Q, K, V, O} do
2:      W\_P $\leftarrow$ weight(A.P)                              Extract
3:      U\_P,  \_P, V\_P^T $\leftarrow$ svd(W\_P)                     Full SVD in FP32
4:      s $\leftarrow$  \_P^2
5:      E $\leftarrow$ cumsum(s) / sum(s)                        Cumulative energy
6:      idx $\leftarrow$ first index where E $\ge$ $\tau$
7:      r\_P $\leftarrow$ max(idx + 1, r\_min)                     Enforce min rank
8:
9:       \_r^{1/2} $\leftarrow$ sqrt( \_P[:r\_P])
10:     U\_r $\leftarrow$ U\_P[:, :r\_P] @ diag( \_r^{1/2})          Symmetric split
11:     V\_r $\leftarrow$ diag( \_r^{1/2}) @ V\_P^T[:r\_P, :]
12:
13:     B\_P $\leftarrow$ zeros(d\_out, r\_L)                       Zero init
14:     A\_P $\leftarrow$ kaiming\_uniform(r\_L, d\_in)
15:
16:     replace A.P with AdaptiveSVDLinear(U\_r, V\_r, A\_P, B\_P)
17: end for
18:
19: freeze(U\_r, V\_r) for all P                        Base frozen
20: unfreeze(A\_P, B\_P) for all P                      Adapters trainable
21:
22: return A'
\end{lstlisting}

\subsection{Complexity Analysis}

\subsubsection{Time Complexity}

\begin{table}[H]
\centering
\small
\begin{tabular}{ll}
\toprule
\textbf{Phase} & \textbf{Time} \\
\midrule
SVD per projection & $\mathcal{O}(d_\text{in} d_\text{out} \min(d_\text{in}, d_\text{out}))$ \\
Rank selection & $\mathcal{O}(\min(d_\text{in}, d_\text{out}))$ \\
Factor truncation & $\mathcal{O}(r_P (d_\text{in} + d_\text{out}))$ \\
Adapter init & $\mathcal{O}(r_L (d_\text{in} + d_\text{out}))$ \\
\textbf{Per projection} & $\mathcal{O}(d_\text{in} d_\text{out} \min(d_\text{in}, d_\text{out}))$ \\
\textbf{Per attention module} & $\mathcal{O}(4 \cdot d^3)$ for square $d \times d$ \\
\textbf{Total ($N_l$ layers)} & $\mathcal{O}(4 N_l d^3)$ \\
\bottomrule
\end{tabular}
\end{table}

For DistilBERT ($N_l = 6$, $d = 768$): ~30 seconds total on GPU.

\subsubsection{Space Complexity}

Peak memory per projection: $O(d^2)$ for the SVD. All covariances 
and calibration states are avoided.

\subsubsection{Compression Time Comparison}

\begin{table}[H]
\centering
\small
\begin{tabular}{llll}
\toprule
\textbf{Method} & \textbf{Calibration} & \textbf{SVD} & \textbf{Total} \\
\midrule
A3 & ~10s & ~20s & ~30s \\
AASVD & ~15s & ~20s & ~35s \\
\textbf{Adaptive SVD} & \textbf{0s} & \textbf{~15s} & \textbf{~15s} \\
\bottomrule
\end{tabular}
\end{table}

\subsection{Parameter Analysis}

For a square $d \times d$ projection:

\begin{table}[H]
\centering
\small
\begin{tabular}{ll}
\toprule
\textbf{Configuration} & \textbf{Parameters} \\
\midrule
Original & $d^2$ \\
Adaptive SVD & $2 r_P d$ \\
Adaptive SVD + LoRA($r_L$) & $2 d (r_P + r_L)$ \\
\bottomrule
\end{tabular}
\end{table}

\textbf{Reduction factor (with adapter):}

\[
\frac{d^2}{2 d (r\_P + r\_L)} = \frac{d}{2(r\_P + r\_L)}
\tag{4.10}
\]

\textbf{Example.} $d = 768$, $r_P = 128$, $r_L = 16$:

\begin{table}[H]
\centering
\small
\begin{tabular}{lll}
\toprule
\textbf{Configuration} & \textbf{Parameters} & \textbf{Ratio} \\
\midrule
Original & 589,824 & 1.00$\times$ \\
Adaptive SVD (base) & 196,608 & 3.00$\times$ \\
+ Adapter & 221,184 & 2.67$\times$ \\
\bottomrule
\end{tabular}
\end{table}

\subsection{Theoretical Notes}

\subsubsection{Reconstruction Bound}

\textbf{Proposition.} With energy threshold $\tau$:

\[
\frac{\|W\_P - U\_{P, r} V\_{P, r}\|\_F}{\|W\_P\|\_F} = \sqrt{1 - \mathcal{E}(r\_P)} \leq \sqrt{1 - \tau}
\tag{4.11}
\]

\textbf{Proof.} $\|W - W_r\|_F^2 = \sum_{i > r} \sigma_i^2 = (1 - \mathcal{E}(r)) \|W\|_F^2$ by 
SVD truncation formula. Since $\mathcal{E}(r_P) \geq \tau$ by rank selection, 
the bound follows. $\square$

At $\tau = 0.92$, the maximum relative error is $\sqrt{0.08} \approx 0.283$.

\subsubsection{Exactness at Full Rank}

\textbf{Proposition.} If $r_P = \mathrm{rank}(W_P)$, then $U_{P, r} V_{P, r} = W_P$ exactly.

\textbf{Proof.} Direct from SVD. $\square$

\subsubsection{Initial Forward Pass Fidelity}

\textbf{Proposition.} At initialization ($B_P = 0$):

\[
\text{Proj}\_P(x) = (x V\_{P, r}^\top) U\_{P, r}^\top + b\_P = W\_{P, r\_P} x + b\_P
\tag{4.12}
\]

\textbf{Proof.} Substitute $B_P = 0$ into (4.9), and note $U_{P, r} V_{P, r} = W_{P, r_P}$. 
$\square$

\textbf{Corollary.} The initial forward pass equals the truncation approximation, 
not the pretrained model (unless $r_P = \text{rank}(W_P)$).

\subsubsection{Rank Reduction Factor}

The parameter ratio per projection is:

\[
\rho\_P = \frac{2 r\_P d}{d^2} = \frac{2 r\_P}{d} \quad \text{(for square } d \times d\text{)}
\tag{4.13}
\]

At $r_P = d/2$, $\rho_P = 1$ (no reduction). At $r_P = d/4$, $\rho_P = 0.5$ (2$\times$ reduction).

\subsection{Hyperparameter Guidance}

\begin{table}[H]
\centering
\small
\begin{tabular}{lll}
\toprule
\textbf{Hyperparameter} & \textbf{Default} & \textbf{Range tested} \\
\midrule
$\tau$ & 0.92 & 0.85 -- 0.99 \\
$r_L$ & 16 & 4 -- 64 \\
$r_\text{min}$ & 16 & 8 -- 32 \\
Recovery steps & 50 & 10 -- 200 \\
\bottomrule
\end{tabular}
\end{table}

\textbf{Choosing $\tau$:}

\begin{itemize}
\item $\tau = 0.85$: aggressive, may lose accuracy
\item $\tau = 0.92$: balanced (default)
\item $\tau = 0.95$: conservative, less compression
\item $\tau = 0.99$: near-lossless, minimal compression

\end{itemize}
\textbf{Choosing $r_L$:}

\begin{itemize}
\item $r_L = 8$: minimal adapter overhead
\item $r_L = 16$: balanced (default)
\item $r_L = 32$: more recovery capacity
\item $r_L = 64$: rarely needed
\end{itemize}